\chapter{Outlook and Conclusion}
\label{ch:bilan}

\section{Technical Summary}
\label{sec:bilan_technique}

The \esibot{} project resulted in a functional differential-drive mobile robot aligned
with recognized architectures such as the TurtleBot3 Burger. Using a standard \ros{}
software stack (\texttt{slam\_toolbox}, Nav2, \texttt{robot\_localization}) to drive a
low-cost robot made it possible to deliver:

\begin{itemize}
	\item a complete, autonomous low-level control chain (encoders, per-wheel \gls{pid},
	      tooled PWM$\to$RPM calibration) entirely running on the Raspberry\,Pi\,4;
	\item a correctly scoped \gls{ekf} sensor fusion (velocities only for odometry,
	      absolute orientation for the \gls{imu}), avoiding double-counting of drift;
	\item a coherent and correctly parameterized Kinect $\to$ \texttt{depthimage\_to\_laserscan}
	      $\to$ \texttt{slam\_toolbox} mapping pipeline (aligned range bounds, loop
	      closure enabled);
	\item a complete Nav2 configuration (\gls{amcl}, dual-layer costmaps, \gls{dwb},
	      recovery behaviors, collision monitor);
	\item functional controller teleoperation and a telemetry web dashboard that are
	      simple to deploy (no external dependency for the HTTP server).
\end{itemize}

\section{Prioritized Action Plan}
\label{sec:bilan_plan}

Based on the analysis in Chapter~\ref{ch:limitations}, the following corrective actions
are proposed in decreasing order of priority:

\begin{enumerate}
	\item \textbf{Reconcile the geometric constants} (wheel radius, wheelbase,
	      \gls{gpio} assignments) between the \gls{urdf}, \texttt{motor\_bridge\_node}
	      and \texttt{motor\_calibrator}, ideally via a single shared parameter file.
	\item \textbf{Conduct the field measurement campaign} described in
	      Chapter~\ref{ch:resultats} (Tables~\ref{tab:test_plan}
	      and~\ref{tab:metrics_consolidated}), in particular the actual mapping of a
	      3\,m$\times$3\,m area with the Kinect, followed by validation of autonomous
	      navigation on that map --- the most urgent priority before the defense, as the
	      jury's demonstration scenario depends directly on it.
	\item \textbf{Implement battery supervision} (\topic{/battery\_state}) to secure
	      trials and demonstrations.
	\item \textbf{Add a minimal vision node} (line or color detection via OpenCV)
	      consuming \topic{/esp32/camera/image\_raw}, even if only for demonstration
	      purposes, to close the gap with Deliverable~3 of the specification.
	\item \textbf{Extend the web dashboard} to include, at minimum, a display of the
	      \gls{slam} map (canvas rendering of the \texttt{OccupancyGrid}) and a
	      low-frequency preview of the camera stream, bringing Deliverable~5 closer to
	      its full specification.
	\item \textbf{Reintroduce an enable button} for controller teleoperation, or
	      otherwise explicitly document this design choice in the README
	      (Deliverable~8).
\end{enumerate}

\section{Future Outlook}
\label{sec:bilan_perspectives}

Beyond the immediate corrections, several avenues for evolution could enrich \esibot{}
in the remainder of the project or for future cohorts: the joint use of the Kinect's
color stream (currently only depth is used) for visual detection complementary to the
ultrasonic radar kept in reserve; the reintroduction of three-sensor fusion (Kinect +
ultrasonic radar + odometry) for increased robustness in the event of a temporary
failure of one source; the addition of a complete power supervision system with
estimation of remaining autonomy; and, more broadly, consolidating the documentation
(README, formal electronic schematic) to make the project more reliably reproducible by
other teams, in keeping with the open-source spirit claimed from the initial
specification onward.

\section{Concluding Remarks}
\label{sec:bilan_conclusion}

This report has documented \esibot{} as it was actually designed and coded, rather than
as it had been initially envisioned on paper --- a deliberate methodological choice
that, the team hopes, constitutes a more useful foundation for the rest of the project
than mere reproduction of the original specification. The identified deviations are not
failures but traceable and justifiable engineering decisions; the action plan above sets
out the path toward complete experimental validation of the system before the final
defense.
